% !TEX root = main.tex
\documentclass[main.tex]{subfiles}
\begin{document}

\section{Derivatives of $\jointdual$}
\label{sec:appendix-joint}

The joint dual $\jointdual(\shift, \weight) = \energy(\dual, \combweight) + \sum_i \weight_i \pvol_i$ defined in \eqref{eq:joint-dual} again depends on $(\shift, \weight)$ only through the shifted sites \eqref{eq:shifted-site}, plus the explicit term $\weight_i \pvol_i$.
Relative to \S\ref{sec:appendix-centroid}, the only new ingredient is the weight reparametrization $\partial_{\weight_i} \combweight_i = 1$ from \eqref{eq:appendix-jac}.
The shift gradient of $\jointdual$ is identical to \eqref{eq:appendix-centroid-grad}, while for the weights the chain rule gives
\begin{equation}
  \partial_{\weight_i} \jointdual
    = (\partial_{\weight_i} \combweight_i) \, (\partial_{\combweight_i} \energy) + \pvol_i
    = (\partial_{\combweight_i} \energy) + \pvol_i
    = -\bigl( \cellvol_i - \pvol_i \bigr) ,
  \label{eq:appendix-joint-grad}
\end{equation}
so that stacking the two gradients recovers the joint gradient \eqref{eq:joint-gradient}.
% 
Applying \eqref{eq:appendix-centroid-hess-chain} a second time, now with $\smash{\partial_{\weight_i} \combweight_i = 1}$, and substituting the energy Hessian \eqref{eq:appendix-energy-hess-off}--\eqref{eq:appendix-energy-hess-diag} gives the off-diagonal blocks
\begin{subequations}
  \label{eq:appendix-blocks-off}
  \begin{empheq}[left={\empheqlbrace\,}]{align}
    \partial^2_{\shift_i \shift_j\tr} \jointdual
      &= \frac{2}{\norm{\dual_i - \dual_j}} \int_{\face_{ij}}\! (\xx - \tcentroid_i)(\xx - \tcentroid_j)\tr \, dA , \label{eq:appendix-blocks-uu} \\
    \partial^2_{\shift_i \weight_j} \jointdual
      &= \frac{1}{\norm{\dual_i - \dual_j}} \int_{\face_{ij}}\! (\xx - \tcentroid_i) \, dA , \label{eq:appendix-blocks-uw} \\
    \partial^2_{\weight_i \weight_j} \jointdual
      &= \frac{1}{2}\frac{\area_{ij}}{\norm{\dual_i - \dual_j}} , \label{eq:appendix-blocks-ww}
  \end{empheq}
\end{subequations}
and diagonal blocks
\begin{subequations}
  \label{eq:appendix-blocks-diag}
  \begin{empheq}[left={\empheqlbrace\,}]{align}
    \partial^2_{\shift_i \shift_i\tr} \jointdual
      &= -\!\sum_{j \sim i} \frac{2}{\norm{\dual_i - \dual_j}} \int_{\face_{ij}}\! (\xx - \tcentroid_i)(\xx - \tcentroid_i)\tr \, dA , \label{eq:appendix-blocks-uu-diag} \\
    \partial^2_{\shift_i \weight_i} \jointdual
      &= -\!\sum_{j \sim i} \frac{1}{\norm{\dual_i - \dual_j}} \int_{\face_{ij}}\! (\xx - \tcentroid_i) \, dA , \label{eq:appendix-blocks-uw-diag} \\
    \partial^2_{\weight_i \weight_i} \jointdual
      &= -\!\sum_{j \sim i} \frac{1}{2}\frac{\area_{ij}}{\norm{\dual_i - \dual_j}} , \label{eq:appendix-blocks-ww-diag}
  \end{empheq}
\end{subequations}
whose shift--shift, shift--weight, and weight--weight parts recover the centroid Hessian \eqref{eq:centroid-hessian}, the mixed block \eqref{eq:joint-hessian}, and the volume Hessian \eqref{eq:newton-step}, respectively.

\end{document}
